% -*- coding: utf-8 -*-
% Aula: Grafos e representação de grafos
% Baseada em: Paulo Feofiloff, "Algoritmos para Grafos (em linguagem C)",
% IME-USP, https://www.ime.usp.br/~pf/algoritmos_para_grafos/
\documentclass[aspectratio=169]{beamer}

\usetheme{UFS}

\usepackage[utf8]{inputenc}
\usepackage[T1]{fontenc}
\usepackage[brazil]{babel}
\usepackage{amsmath,amssymb,amsthm}
\usepackage{graphicx}
\usepackage{booktabs}
\usepackage{listings}
\usepackage{xcolor}
\usepackage{tikz}
\usetikzlibrary{arrows.meta, shapes, positioning, calc, fit, matrix,
  decorations.pathreplacing, backgrounds, chains, shapes.geometric}

\definecolor{codegreen}{rgb}{0,0.6,0}
\definecolor{codegray}{rgb}{0.5,0.5,0.5}
\definecolor{codepurple}{rgb}{0.58,0,0.82}
\definecolor{backcolour}{rgb}{0.95,0.95,0.92}

\lstdefinestyle{mystyle}{
    language=C,
    backgroundcolor=\color{backcolour},
    commentstyle=\color{codegreen},
    keywordstyle=\color{magenta},
    numberstyle=\tiny\color{codegray},
    stringstyle=\color{codepurple},
    breakatwhitespace=false,
    breaklines=true,
    captionpos=b,
    keepspaces=true,
    numbers=none,
    showspaces=false,
    showstringspaces=false,
    showtabs=false,
    tabsize=2,
    frame=single,
    rulecolor=\color{gray},
    extendedchars=true,
    basicstyle=\ttfamily\scriptsize,
    literate={ê}{{\^e}}1 {à}{{\`a}}1 {ú}{{\'u}}1 {õ}{{\~o}}1 {ó}{{\'o}}1
             {í}{{\'i}}1 {á}{{\'a}}1 {ã}{{\~a}}1 {é}{{\'e}}1 {ç}{{\c{c}}}1
             {â}{{\^a}}1 {ô}{{\^o}}1 {ü}{{\"u}}1 {Á}{{\'A}}1 {É}{{\'E}}1
}
\lstset{style=mystyle}

% ---- estilos TikZ reutilizados nos diagramas de grafos ----
\tikzset{
  vtx/.style={circle, draw=blue!70, fill=blue!10, minimum size=6mm,
              inner sep=0pt, font=\scriptsize},
  vtxhl/.style={circle, draw=orange!80!black, fill=orange!25, minimum size=6mm,
                inner sep=0pt, font=\scriptsize},
  arc/.style={-{Stealth[length=2mm]}, thick, blue!60!black},
  arcgray/.style={-{Stealth[length=2mm]}, thick, gray!60},
  edg/.style={thick, blue!60!black},
  cell/.style={draw=gray!60, minimum size=4.2mm, inner sep=0pt, font=\tiny},
  lnode/.style={draw=blue!60, fill=blue!8, minimum height=4.5mm,
                minimum width=4.5mm, inner sep=1pt, font=\tiny},
  box/.style={rectangle, draw=blue!60, fill=blue!10, rounded corners,
              minimum width=4.6cm, minimum height=0.55cm, font=\scriptsize}
}

\title{Grafos e representação de grafos}
\subtitle{Algoritmos e Estruturas de Dados 2}
\author[André Y. Kashiwabara <andreyoshiaki@dcomp.ufs.br>]{Prof.\ Dr.\ André Yoshiaki Kashiwabara}
\institute{Departamento de Computação --- Universidade Federal de Sergipe}
\date{\today}

\begin{document}

\begin{frame}[plain,noframenumbering]
  \titlepage
\end{frame}

% =====================================================================
\begin{frame}{O que você vai aprender hoje}
  \begin{center}
  \begin{tikzpicture}[node distance=0.18cm]
    \node[box] (t1) at (0,0)     {1. O que é um grafo: definições básicas};
    \node[box] (t2) at (0,-0.78) {2. Leques, graus e contagem de arcos};
    \node[box] (t3) at (0,-1.56) {3. Grafos não-dirigidos e subgrafos};
    \node[box] (t4) at (0,-2.34) {4. Matriz de adjacências (em C)};
    \node[box] (t5) at (0,-3.12) {5. Listas de adjacência (em C)};
    \node[box] (t6) at (0,-3.90) {6. Qual representação escolher?};
    \draw[arc] (t1)--(t2);
    \draw[arc] (t2)--(t3);
    \draw[arc] (t3)--(t4);
    \draw[arc] (t4)--(t5);
    \draw[arc] (t5)--(t6);
  \end{tikzpicture}
  \end{center}
  \vfill
  {\scriptsize Material baseado em Feofiloff, \emph{Algoritmos para Grafos (em linguagem C)}, IME-USP.}
\end{frame}

% =====================================================================
\section{Definições básicas}

\begin{frame}{Por que precisamos de grafos?}
  \begin{columns}[T]
    \begin{column}{0.48\textwidth}
      \textbf{Problema}

      \medskip
      {\small
      Muitos problemas são sobre \emph{objetos} e \emph{ligações} entre eles:
      \begin{itemize}\scriptsize
        \item páginas da Web e os \emph{links} entre elas;
        \item cidades e estradas;
        \item tarefas e suas dependências;
        \item pessoas e amizades.
      \end{itemize}
      Precisamos de um único modelo para todos eles.}
    \end{column}
    \begin{column}{0.48\textwidth}
      \textbf{Solução}

      \medskip
      {\small
      O \textbf{grafo}: um modelo combinatório mínimo --- um conjunto de
      \emph{vértices} e um conjunto de \emph{arcos}.

      \medskip
      Resolvido o problema no grafo, a solução vale para todas as aplicações
      que ele modela.}
    \end{column}
  \end{columns}
\end{frame}

\begin{frame}{Grafo: definição}
  \begin{block}{Definição}
    {\small Um \textbf{grafo} é um par de conjuntos: um conjunto de
    \textbf{vértices} e um conjunto de \textbf{arcos}. Cada \textbf{arco} é um
    \emph{par ordenado} de vértices: o primeiro é a \textbf{ponta inicial} e o
    segundo é a \textbf{ponta final}. O arco de ponta inicial $v$ e ponta final
    $w$ é denotado por $v$\texttt{-}$w$.}
  \end{block}
  \begin{exampleblock}{Intuição e convenções deste curso}
    {\small
    \begin{itemize}\scriptsize
      \item Arcos são \emph{dirigidos}: $v$\texttt{-}$w$ sai de $v$ e entra em $w$.
      \item A presença de $v$\texttt{-}$w$ é independente da presença de $w$\texttt{-}$v$.
      \item \textbf{Não há laços}: a ponta final é sempre diferente da inicial.
      \item \textbf{Não há arcos paralelos}: arcos são \emph{meros pares} de vértices,
            logo dois arcos distintos não podem ter as mesmas pontas.
    \end{itemize}}
  \end{exampleblock}
  {\scriptsize ``Grafo dirigido'' (ou \emph{digrafo}) é sinônimo; o adjetivo é redundante aqui.}
\end{frame}

\begin{frame}{Um grafo concreto}
  {\small Grafo sobre os vértices $0..5$ com o conjunto de arcos}
  \begin{center}\ttfamily\small 0-1\quad 0-5\quad 1-0\quad 1-5\quad 2-4\quad 3-1\quad 5-3\end{center}
  \begin{center}
  \begin{tikzpicture}[scale=1.0]
    \node[vtx] (v0) at (0,1.4)  {0};
    \node[vtx] (v1) at (2.2,1.4){1};
    \node[vtx] (v2) at (4.4,1.4){2};
    \node[vtx] (v3) at (0,0)    {3};
    \node[vtx] (v5) at (2.2,0)  {5};
    \node[vtx] (v4) at (4.4,0)  {4};
    \draw[arc] (v0) to[bend left=18]  (v1);
    \draw[arc] (v1) to[bend left=18]  (v0);
    \draw[arc] (v0) -- (v5);
    \draw[arc] (v1) -- (v5);
    \draw[arc] (v2) -- (v4);
    \draw[arc] (v3) -- (v1);
    \draw[arc] (v5) -- (v3);
  \end{tikzpicture}
  \end{center}
  {\scriptsize $0$\texttt{-}$1$ e $1$\texttt{-}$0$ são \textbf{antiparalelos}: a ponta inicial de um é a ponta final do outro.
  Note que $4$ não tem nenhum arco saindo dele.}
\end{frame}

\begin{frame}{Vizinhança, adjacência e tamanho}
  \begin{block}{Definição}
    {\small Um vértice $w$ é \textbf{vizinho} de $v$ (ou \textbf{adjacente} a $v$)
    se $v$\texttt{-}$w$ é um arco do grafo.

    O \textbf{tamanho} de um grafo com $V$ vértices e $A$ arcos é a soma $V+A$.}
  \end{block}
  \begin{alertblock}{Cuidado}
    {\small A relação de vizinhança \textbf{não é simétrica}: $w$ pode ser
    vizinho de $v$ sem que $v$ seja vizinho de $w$.}
  \end{alertblock}
  {\small No grafo anterior: os vizinhos de $0$ são $1$ e $5$; o vizinho de $3$ é $1$;
  o vértice $4$ não tem vizinho algum. Ali $V=6$, $A=7$ e o tamanho é $13$.}
\end{frame}

\begin{frame}{Exemplo: a Web é um grafo}
  \begin{exampleblock}{Exemplo}
    {\small A estrutura da rede WWW pode ser representada por um grafo: os
    \textbf{vértices} são as páginas HTML e os \textbf{arcos} são os \emph{links}
    que apontam de uma página para outra. Navegar na rede é pular de um vértice
    a outro seguindo os arcos.}
  \end{exampleblock}
  \begin{center}
  \begin{tikzpicture}[scale=0.9]
    \node[vtx] (a) at (0,0.9)   {A};
    \node[vtx] (b) at (2.0,1.5) {B};
    \node[vtx] (c) at (2.0,0.2) {C};
    \node[vtx] (d) at (4.0,0.9) {D};
    \draw[arc] (a) -- (b);
    \draw[arc] (a) -- (c);
    \draw[arc] (b) -- (d);
    \draw[arc] (c) -- (d);
    \draw[arc] (d) to[bend left=35] (a);
  \end{tikzpicture}
  \end{center}
  {\scriptsize A direção importa: A aponta para B, mas B não aponta para A.
  É exatamente isso que a definição de arco como \emph{par ordenado} captura.}
\end{frame}

\begin{frame}{Recapitulando: definições básicas}
  \begin{block}{O que aprendemos}
    \begin{itemize}\small
      \item Grafo $=$ conjunto de vértices $+$ conjunto de arcos (pares ordenados).
      \item Notação $v$\texttt{-}$w$: sai de $v$, entra em $w$.
      \item $w$ é vizinho (adjacente) de $v$ se $v$\texttt{-}$w$ é arco; relação não simétrica.
      \item Sem laços e sem arcos paralelos, por convenção deste curso.
      \item $V$ vértices, $A$ arcos, tamanho $V+A$.
    \end{itemize}
  \end{block}
\end{frame}

% =====================================================================
\section{Leques e graus}

\begin{frame}{Por que precisamos de graus?}
  \begin{columns}[T]
    \begin{column}{0.48\textwidth}
      \textbf{Problema}

      \medskip
      {\small Como medir a ``importância local'' de um vértice? Como saber
      quantos arcos um grafo pode ter, e portanto quanta memória vamos gastar
      para guardá-lo?}
    \end{column}
    \begin{column}{0.48\textwidth}
      \textbf{Solução}

      \medskip
      {\small Contar os arcos que saem e que entram em cada vértice: os
      \textbf{graus}. Com eles obtemos limites sobre $A$ e classificamos grafos
      em \emph{densos} e \emph{esparsos}.}
    \end{column}
  \end{columns}
\end{frame}

\begin{frame}{Leques e graus: definição}
  \begin{block}{Definição}
    {\small
    O \textbf{leque de saída} (\emph{fan-out}) de um vértice é o conjunto de
    todos os arcos que saem dele; o \textbf{leque de entrada} (\emph{fan-in}) é
    o conjunto dos arcos que entram nele.

    \smallskip
    O \textbf{grau de saída} (\emph{outdegree}) é o tamanho do leque de saída;
    o \textbf{grau de entrada} (\emph{indegree}) é o tamanho do leque de entrada.}
  \end{block}
  \begin{exampleblock}{Vértices com nome próprio}
    {\small
    \begin{itemize}\scriptsize
      \item \textbf{fonte} (\emph{source}): grau de entrada nulo;
      \item \textbf{sorvedouro} (\emph{sink}): grau de saída nulo;
      \item \textbf{isolado}: graus de entrada e de saída ambos nulos.
    \end{itemize}}
  \end{exampleblock}
\end{frame}

\begin{frame}{Leques, visualmente}
  \begin{center}
  \begin{tikzpicture}[scale=1.0]
    \node[vtxhl] (v) at (2.4,0.8) {$v$};
    \node[vtx] (a) at (0,1.6) {a};
    \node[vtx] (b) at (0,0.0) {b};
    \node[vtx] (c) at (4.8,1.9) {c};
    \node[vtx] (d) at (4.8,0.8) {d};
    \node[vtx] (e) at (4.8,-0.3) {e};
    \draw[arc] (a) -- (v);
    \draw[arc] (b) -- (v);
    \draw[arc] (v) -- (c);
    \draw[arc] (v) -- (d);
    \draw[arc] (v) -- (e);
    \node[font=\scriptsize, gray!70!black] at (0.5,-0.9) {leque de entrada: 2 arcos};
    \node[font=\scriptsize, gray!70!black] at (4.6,-1.1) {leque de saída: 3 arcos};
  \end{tikzpicture}
  \end{center}
  {\small Aqui o grau de entrada de $v$ é $2$ e o grau de saída é $3$.
  Se o leque de entrada fosse vazio, $v$ seria uma \textbf{fonte};
  se o leque de saída fosse vazio, um \textbf{sorvedouro}.}
\end{frame}

\begin{frame}{Quantos arcos um grafo pode ter?}
  \begin{block}{Propriedade}
    {\small A soma dos graus de \emph{saída} de todos os vértices é igual ao
    número de arcos do grafo. O mesmo vale para a soma dos graus de
    \emph{entrada}.}
  \end{block}
  \begin{alertblock}{Consequência}
    {\small Um grafo com $V$ vértices tem no máximo $V(V-1)$ arcos ---
    um pouco menos que $V^2$.}
  \end{alertblock}
  {\small
  \begin{itemize}\small
    \item Um grafo é \textbf{completo} se todo par ordenado de vértices
          distintos é um arco: tem exatamente $V(V-1)$ arcos.
    \item Um \textbf{torneio} tem, para cada par $v\,w$ de vértices distintos,
          exatamente um dentre $v$\texttt{-}$w$ e $w$\texttt{-}$v$:
          tem $\tfrac{1}{2}V(V-1)$ arcos.
  \end{itemize}}
\end{frame}

\begin{frame}{Densos e esparsos}
  \begin{columns}[T]
    \begin{column}{0.5\textwidth}
      {\small
      \textbf{Denso}: $A$ é da mesma ordem que $V^2$
      (digamos $V^2/2$, $V^2/10$, \dots). O tamanho é proporcional a $V^2$.

      \medskip
      \textbf{Esparso}: $A$ é da mesma ordem que $V$
      (digamos $10V$, $V/2$, \dots). O tamanho é proporcional a $V$.

      \medskip
      {\scriptsize Estas definições se aplicam a \emph{famílias infinitas} de
      grafos, não a um grafo individual.}}
    \end{column}
    \begin{column}{0.5\textwidth}
      \begin{center}
      \begin{tikzpicture}[scale=0.75]
        \draw[->] (0,0) -- (4.3,0) node[right,font=\tiny] {$V$};
        \draw[->] (0,0) -- (0,3.2) node[above,font=\tiny] {$A$};
        \draw[thick,red!70!black,domain=0:1.75,smooth,variable=\x]
          plot ({\x*2.28},{\x*\x});
        \node[red!70!black,font=\tiny,right] at (3.6,2.6) {denso $\sim V^2$};
        \draw[thick,blue!70!black] (0,0) -- (4.0,0.85);
        \node[blue!70!black,font=\tiny,right] at (3.0,0.45) {esparso $\sim V$};
      \end{tikzpicture}
      \end{center}
      {\scriptsize A diferença decide qual estrutura de dados vale a pena.}
    \end{column}
  \end{columns}
\end{frame}

\begin{frame}{Recapitulando: leques e graus}
  \begin{block}{O que aprendemos}
    \begin{itemize}\small
      \item Leque de saída/entrada; graus de saída (\emph{outdegree}) e de entrada (\emph{indegree}).
      \item Fonte, sorvedouro e vértice isolado.
      \item Soma dos graus de saída $=$ soma dos graus de entrada $= A$.
      \item No máximo $V(V-1)$ arcos; completo tem $V(V-1)$, torneio tem $\tfrac{1}{2}V(V-1)$.
      \item Denso ($A \sim V^2$) versus esparso ($A \sim V$).
    \end{itemize}
  \end{block}
\end{frame}

% =====================================================================
\section{Grafos não-dirigidos e subgrafos}

\begin{frame}{Por que precisamos de grafos não-dirigidos?}
  \begin{columns}[T]
    \begin{column}{0.48\textwidth}
      \textbf{Problema}

      \medskip
      {\small Amizade, estradas de mão dupla e cabos de rede são relações
      \emph{simétricas}. Modelá-las com arcos dirigidos obriga a lembrar sempre
      de inserir os dois sentidos.}
    \end{column}
    \begin{column}{0.48\textwidth}
      \textbf{Solução}

      \medskip
      {\small Tratar cada par de arcos antiparalelos como uma entidade nova ---
      a \textbf{aresta}. O grafo continua sendo o mesmo objeto; muda apenas
      como falamos dele.}
    \end{column}
  \end{columns}
\end{frame}

\begin{frame}{Grafo não-dirigido: definição}
  \begin{block}{Definição}
    {\small Um grafo é \textbf{não-dirigido} se cada um de seus arcos é
    antiparalelo a algum outro arco: para todo arco $v$\texttt{-}$w$, o grafo
    também tem o arco $w$\texttt{-}$v$.

    \smallskip
    Cada par de arcos antiparalelos é uma \textbf{aresta} (\emph{edge}).
    Se $E$ é o número de arestas e $A$ o de arcos, então $E = A/2$.}
  \end{block}
  \begin{exampleblock}{Vocabulário}
    {\small Uma aresta $v$\texttt{-}$w$ \textbf{liga} $v$ e $w$, e
    \textbf{incide} em $v$ e em $w$. O \textbf{leque} de $v$ é o conjunto de
    arestas que incidem em $v$; o \textbf{grau} de $v$ é o número dessas
    arestas. A soma dos graus vale $2E$.}
  \end{exampleblock}
\end{frame}

\begin{frame}{Um grafo não-dirigido}
  {\small O conjunto de \emph{arcos}}
  \begin{center}\ttfamily\scriptsize
    0-1 0-2 0-5 1-2 2-3 2-4 3-4 3-5\quad
    1-0 2-0 5-0 2-1 3-2 4-2 4-3 5-3
  \end{center}
  {\small define um grafo não-dirigido, cujo conjunto de \emph{arestas} é
  \texttt{0-1 0-2 0-5 1-2 2-3 2-4 3-4 3-5} ($A=16$, $E=8$).}
  \begin{center}
  \begin{tikzpicture}[scale=1.0]
    \node[vtx] (v0) at (0,1.3)   {0};
    \node[vtx] (v1) at (1.6,2.1) {1};
    \node[vtx] (v2) at (3.0,1.3) {2};
    \node[vtx] (v3) at (4.4,0.6) {3};
    \node[vtx] (v4) at (3.0,-0.2){4};
    \node[vtx] (v5) at (1.4,0.0) {5};
    \draw[edg] (v0)--(v1); \draw[edg] (v0)--(v2); \draw[edg] (v0)--(v5);
    \draw[edg] (v1)--(v2); \draw[edg] (v2)--(v3); \draw[edg] (v2)--(v4);
    \draw[edg] (v3)--(v4); \draw[edg] (v3)--(v5);
  \end{tikzpicture}
  \end{center}
  {\scriptsize A adjacência aqui é simétrica; o grau de $2$ é $4$ e o de $1$ é $2$.}
\end{frame}

\begin{frame}{Subgrafos}
  \begin{block}{Definição}
    {\small Um grafo $H$ é \textbf{subgrafo} de um grafo $G$ (notação
    $H \subseteq G$) se todo vértice de $H$ é vértice de $G$ e todo arco de $H$
    é arco de $G$.}
  \end{block}
  {\small
  \begin{itemize}\small
    \item \textbf{gerador} (\emph{spanning}): contém \emph{todos} os vértices de $G$;
    \item \textbf{induzido}: todo arco de $G$ com ambas as pontas em $H$ também é arco de $H$;
    \item \textbf{próprio}: diferente de $G$ (notação $H \subset G$).
  \end{itemize}

  \smallskip
  Dado um conjunto $X$ de vértices, $G[X]$ é o subgrafo induzido por $X$.
  Se $a$ é um arco, $G-a$ é o subgrafo obtido removendo $a$.}
\end{frame}

\begin{frame}{Exemplo: subgrafo induzido}
  \begin{center}
  \begin{tikzpicture}[scale=0.95]
    \begin{scope}
      \node[font=\scriptsize] at (1.6,2.6) {$G$};
      \node[vtx] (a0) at (0,1.3)   {0};
      \node[vtx] (a1) at (1.6,2.1) {1};
      \node[vtx] (a2) at (3.0,1.3) {2};
      \node[vtx] (a3) at (4.0,0.3) {3};
      \node[vtx] (a4) at (2.6,-0.4){4};
      \node[vtx] (a5) at (1.0,-0.2){5};
      \draw[edg] (a0)--(a1); \draw[edg] (a0)--(a2); \draw[edg] (a0)--(a5);
      \draw[edg] (a1)--(a2); \draw[edg] (a2)--(a3); \draw[edg] (a2)--(a4);
      \draw[edg] (a3)--(a4); \draw[edg] (a3)--(a5);
    \end{scope}
    \begin{scope}[xshift=6.0cm]
      \node[font=\scriptsize] at (1.6,2.6) {$G[\{0,1,2,4\}]$};
      \node[vtxhl] (b0) at (0,1.3)   {0};
      \node[vtxhl] (b1) at (1.6,2.1) {1};
      \node[vtxhl] (b2) at (3.0,1.3) {2};
      \node[vtxhl] (b4) at (2.6,-0.4){4};
      \draw[edg] (b0)--(b1); \draw[edg] (b0)--(b2);
      \draw[edg] (b1)--(b2); \draw[edg] (b2)--(b4);
    \end{scope}
  \end{tikzpicture}
  \end{center}
  {\small O subgrafo induzido por $X=\{0,1,2,4\}$ mantém \emph{todas} as arestas
  de $G$ com as duas pontas em $X$ --- nem uma a mais, nem uma a menos.}
\end{frame}

\begin{frame}{Recapitulando: não-dirigidos e subgrafos}
  \begin{block}{O que aprendemos}
    \begin{itemize}\small
      \item Não-dirigido: todo arco tem seu antiparalelo; cada par é uma \textbf{aresta}.
      \item $E = A/2$; soma dos graus $= 2E$.
      \item Subgrafo $H \subseteq G$; gerador, induzido, próprio.
      \item $G[X]$: subgrafo induzido por um conjunto $X$ de vértices.
      \item Um subgrafo de um grafo não-dirigido pode \emph{não} ser não-dirigido.
    \end{itemize}
  \end{block}
\end{frame}

% =====================================================================
\section{Matriz de adjacências}

\begin{frame}[fragile]{Por que precisamos de uma estrutura de dados?}
  \begin{columns}[T]
    \begin{column}{0.48\textwidth}
      \textbf{Problema}

      \medskip
      {\small A definição matemática de grafo não diz \emph{nada} sobre como
      guardá-lo na memória. Sem uma representação concreta não há algoritmo,
      não há programa em C, não há análise de consumo de tempo.}
    \end{column}
    \begin{column}{0.48\textwidth}
      \textbf{Solução}

      \medskip
      {\small Duas representações clássicas. Antes delas, uma convenção:

      \smallskip
      \lstinline|#define vertex int|

      \smallskip
      Os vértices são sempre $0\;1\;2\;\ldots\;V-1$.}
    \end{column}
  \end{columns}
\end{frame}

\begin{frame}{Matriz de adjacências: definição}
  \begin{block}{Definição}
    {\small A \textbf{matriz de adjacências} de um grafo é uma matriz booleana
    com linhas e colunas indexadas pelos vértices:
    \[
      \texttt{adj[v][w]} =
      \begin{cases}
        1 & \text{se } v\text{\texttt{-}}w \text{ é um arco},\\
        0 & \text{caso contrário.}
      \end{cases}
    \]}
  \end{block}
  \begin{exampleblock}{Leitura da matriz}
    {\small A \emph{linha} $v$ representa o leque de \textbf{saída} de $v$;
    a \emph{coluna} $w$ representa o leque de \textbf{entrada} de $w$.
    Como não há laços, a diagonal é toda nula. Se o grafo é não-dirigido,
    a matriz é \emph{simétrica}: \texttt{adj[v][w] $\equiv$ adj[w][v]}.}
  \end{exampleblock}
\end{frame}

\begin{frame}{A matriz do nosso exemplo}
  {\small Arcos \ttfamily 0-1 0-5 1-0 1-5 2-4 3-1 5-3\rmfamily :}
  \begin{center}
  \begin{tikzpicture}[scale=0.95]
    \begin{scope}
      \node[vtx] (v0) at (0,1.4)  {0};
      \node[vtx] (v1) at (1.7,1.4){1};
      \node[vtx] (v2) at (3.4,1.4){2};
      \node[vtx] (v3) at (0,0)    {3};
      \node[vtx] (v5) at (1.7,0)  {5};
      \node[vtx] (v4) at (3.4,0)  {4};
      \draw[arc] (v0) to[bend left=18] (v1);
      \draw[arc] (v1) to[bend left=18] (v0);
      \draw[arc] (v0) -- (v5); \draw[arc] (v1) -- (v5);
      \draw[arc] (v2) -- (v4); \draw[arc] (v3) -- (v1);
      \draw[arc] (v5) -- (v3);
    \end{scope}
    \begin{scope}[xshift=5.4cm, yshift=-0.3cm]
      \foreach \j/\lab in {0/0,1/1,2/2,3/3,4/4,5/5}
        \node[font=\tiny,gray!60!black] at (\j*0.46+0.46,2.05) {\lab};
      \foreach \i/\lab in {0/0,1/1,2/2,3/3,4/4,5/5}
        \node[font=\tiny,gray!60!black] at (0,1.75-\i*0.46) {\lab};
      \foreach \i/\row in {0/{0,1,0,0,0,1}, 1/{1,0,0,0,0,1}, 2/{0,0,0,0,1,0},
                           3/{0,1,0,0,0,0}, 4/{0,0,0,0,0,0}, 5/{0,0,0,1,0,0}} {
        \foreach \val [count=\j from 0] in \row {
          \node[cell, fill={\ifnum\val=1 orange!30\else white\fi}]
            at (\j*0.46+0.46,1.75-\i*0.46) {\val};
        }
      }
    \end{scope}
  \end{tikzpicture}
  \end{center}
  {\scriptsize A linha $4$ é toda nula: $4$ é um \textbf{sorvedouro}.
  A coluna $2$ é toda nula: $2$ é uma \textbf{fonte}.}
\end{frame}

\begin{frame}[fragile]{Matriz: a estrutura em C}
\begin{lstlisting}
/* REPRESENTACAO POR MATRIZ DE ADJACENCIAS:
   O campo adj e um ponteiro para a matriz de adjacencias do grafo.
   O campo V contem o numero de vertices e A o numero de arcos. */
struct graph {
   int V;
   int A;
   int **adj;
};

/* Um Graph e um ponteiro para um graph. */
typedef struct graph *Graph;
\end{lstlisting}
  {\small O grafo carrega consigo $V$ e $A$: assim, perguntar quantos vértices
  ou quantos arcos o grafo tem custa tempo constante.}
\end{frame}

\begin{frame}[fragile]{Matriz: construção do grafo}
\begin{lstlisting}[basicstyle=\ttfamily\tiny]
/* Constroi um grafo com vertices 0 1 .. V-1 e nenhum arco. */
Graph GRAPHinit( int V) {
   Graph G = malloc( sizeof *G);
   G->V = V;
   G->A = 0;
   G->adj = MATRIXint( V, V, 0);
   return G;
}

/* Aloca matriz com linhas 0..r-1 e colunas 0..c-1, tudo com valor val. */
static int **MATRIXint( int r, int c, int val) {
   int **m = malloc( r * sizeof (int *));
   for (vertex i = 0; i < r; ++i)
      m[i] = malloc( c * sizeof (int));
   for (vertex i = 0; i < r; ++i)
      for (vertex j = 0; j < c; ++j)
         m[i][j] = val;
   return m;
}
\end{lstlisting}
\end{frame}

\begin{frame}[fragile]{Matriz: inserção e remoção de arcos}
\begin{lstlisting}[basicstyle=\ttfamily\tiny]
/* Insere um arco v-w no grafo G. Supoe v e w distintos, positivos
   e menores que G->V. Se o arco ja existe, nao faz nada. */
void GRAPHinsertArc( Graph G, vertex v, vertex w) {
   if (G->adj[v][w] == 0) {
      G->adj[v][w] = 1;
      G->A++;
   }
}

/* Remove do grafo G o arco v-w. Se o arco nao existe, nao faz nada. */
void GRAPHremoveArc( Graph G, vertex v, vertex w) {
   if (G->adj[v][w] == 1) {
      G->adj[v][w] = 0;
      G->A--;
   }
}
\end{lstlisting}
  {\small Ambas consomem tempo \textbf{constante}: um acesso à matriz decide tudo.}
\end{frame}

\begin{frame}[fragile]{Matriz: listando os vizinhos}
\begin{lstlisting}
/* Imprime, para cada vertice v do grafo G, em uma linha,
   todos os vertices adjacentes a v. */
void GRAPHshow( Graph G) {
   for (vertex v = 0; v < G->V; ++v) {
      printf( "%2d:", v);
      for (vertex w = 0; w < G->V; ++w)
         if (G->adj[v][w] == 1)
            printf( " %2d", w);
      printf( "\n");
   }
}
\end{lstlisting}
  \begin{alertblock}{Custo}
    {\scriptsize Percorrer os vizinhos de \emph{um} vértice custa tempo
    proporcional a $V$, mesmo que ele tenha um único vizinho:
    \texttt{GRAPHshow()} custa $V^2$.}
  \end{alertblock}
\end{frame}

\begin{frame}{Recapitulando: matriz de adjacências}
  \begin{block}{O que aprendemos}
    \begin{itemize}\small
      \item \texttt{adj[v][w]} $\in \{0,1\}$; linha $=$ leque de saída, coluna $=$ leque de entrada.
      \item Diagonal nula (sem laços); matriz simétrica se o grafo é não-dirigido.
      \item Espaço proporcional a $V^2$ --- \emph{proporcional ao tamanho} apenas se o grafo é denso.
      \item Testar se $v$\texttt{-}$w$ é arco, inserir e remover: tempo constante.
      \item Percorrer os vizinhos de um vértice: tempo proporcional a $V$.
    \end{itemize}
  \end{block}
\end{frame}

% =====================================================================
\section{Listas de adjacência}

\begin{frame}{Por que precisamos de listas de adjacência?}
  \begin{columns}[T]
    \begin{column}{0.48\textwidth}
      \textbf{Problema}

      \medskip
      {\small Um grafo esparso com $V = 10^6$ vértices e $A = 5 \cdot 10^6$ arcos
      exigiria $10^{12}$ células de matriz --- quase todas iguais a zero.
      Impossível na prática.}
    \end{column}
    \begin{column}{0.48\textwidth}
      \textbf{Solução}

      \medskip
      {\small Guardar, para cada vértice, apenas a lista dos vizinhos que ele
      realmente tem. O espaço passa a ser proporcional ao \textbf{tamanho}
      $V+A$ do grafo.}
    \end{column}
  \end{columns}
\end{frame}

\begin{frame}{Listas de adjacência: definição}
  \begin{block}{Definição}
    {\small O \textbf{vetor de listas de adjacência} de um grafo tem uma lista
    encadeada associada a cada vértice. A lista associada a $v$ contém todos os
    vizinhos de $v$; portanto, ela representa o \textbf{leque de saída} de $v$.}
  \end{block}
  \begin{exampleblock}{Exemplo}
    {\small Para os arcos \texttt{0-1 0-5 1-0 1-5 2-4 3-1 5-3}:}
    \begin{center}\ttfamily\scriptsize
      0: 5 1 \quad\quad 1: 5 0 \quad\quad 2: 4 \quad\quad
      3: 1 \quad\quad 4: \quad\quad 5: 3
    \end{center}
    {\scriptsize A \emph{ordem} dentro de cada lista depende da ordem de
    inserção --- ela não faz parte da definição do grafo.}
  \end{exampleblock}
\end{frame}

\begin{frame}{O vetor de listas, visualmente}
  \begin{center}
  \begin{tikzpicture}[scale=1.0]
    \foreach \i/\lab in {0/0,1/1,2/2,3/3,4/4,5/5} {
      \node[lnode, fill=blue!18] (a\i) at (0,-\i*0.62) {\lab};
    }
    \node[font=\tiny, gray!60!black, above] at (0,0.42) {\texttt{adj}};
    % linha 0: 5 -> 1
    \node[lnode] (n01) at (1.5,0)      {5};  \node[lnode] (n02) at (2.7,0)      {1};
    \node[lnode] (n11) at (1.5,-0.62)  {5};  \node[lnode] (n12) at (2.7,-0.62)  {0};
    \node[lnode] (n21) at (1.5,-1.24)  {4};
    \node[lnode] (n31) at (1.5,-1.86)  {1};
    \node[font=\tiny] (n41) at (1.5,-2.48) {\texttt{NULL}};
    \node[lnode] (n51) at (1.5,-3.10)  {3};
    \foreach \i/\dst in {0/n01, 1/n11, 2/n21, 3/n31, 4/n41, 5/n51}
      \draw[arcgray] (a\i) -- (\dst);
    \draw[arcgray] (n01) -- (n02);
    \draw[arcgray] (n11) -- (n12);
    \node[font=\tiny] at (3.9,0)     {\texttt{NULL}};
    \node[font=\tiny] at (3.9,-0.62) {\texttt{NULL}};
    \node[font=\tiny] at (2.7,-1.24) {\texttt{NULL}};
    \node[font=\tiny] at (2.7,-1.86) {\texttt{NULL}};
    \node[font=\tiny] at (2.7,-3.10) {\texttt{NULL}};
  \end{tikzpicture}
  \end{center}
  {\scriptsize O vetor tem $V$ posições; somando todos os nós das listas temos
  exatamente $A$ nós. Daí o espaço proporcional a $V+A$.}
\end{frame}

\begin{frame}[fragile]{Listas: as estruturas em C}
\begin{lstlisting}
/* Cada no da lista corresponde a um arco: contem um vizinho w de v
   e o endereco do no seguinte. Um link e um ponteiro para um node. */
typedef struct node *link;
struct node {
   vertex w;
   link next;
};

/* O campo adj e um ponteiro para o vetor de listas de adjacencia. */
struct graph {
   int V;
   int A;
   link *adj;
};
typedef struct graph *Graph;
\end{lstlisting}
\end{frame}

\begin{frame}[fragile]{Listas: construção do grafo}
\begin{lstlisting}[basicstyle=\ttfamily\tiny]
/* Recebe um vertice w e o endereco next de um no e devolve o endereco
   a de um novo no tal que a->w == w e a->next == next. */
static link NEWnode( vertex w, link next) {
   link a = malloc( sizeof (struct node));
   a->w = w;
   a->next = next;
   return a;
}

/* Constroi um grafo com vertices 0 1 .. V-1 e nenhum arco. */
Graph GRAPHinit( int V) {
   Graph G = malloc( sizeof *G);
   G->V = V;
   G->A = 0;
   G->adj = malloc( V * sizeof (link));
   for (vertex v = 0; v < V; ++v)
      G->adj[v] = NULL;
   return G;
}
\end{lstlisting}
\end{frame}

\begin{frame}[fragile]{Listas: inserção de arco}
\begin{lstlisting}
/* Insere um arco v-w no grafo G. Supoe v e w distintos, positivos
   e menores que G->V. Se o grafo ja tem o arco v-w, nao faz nada. */
void GRAPHinsertArc( Graph G, vertex v, vertex w) {
   for (link a = G->adj[v]; a != NULL; a = a->next)
      if (a->w == w) return;
   G->adj[v] = NEWnode( w, G->adj[v]);
   G->A++;
}
\end{lstlisting}
  \begin{alertblock}{Custo}
    {\small A função consome muito tempo (no pior caso, tempo proporcional ao
    \emph{número de arcos}), pois percorre a lista para verificar se o arco a
    inserir já existe. O novo nó entra sempre no \emph{início} da lista --- por
    isso a ordem observada no exemplo é o inverso da ordem de inserção.}
  \end{alertblock}
\end{frame}

\begin{frame}[fragile]{Recapitulando: listas de adjacência}
  \begin{block}{O que aprendemos}
    \begin{itemize}\small
      \item Vetor de $V$ listas encadeadas; a lista de $v$ é o leque de saída de $v$.
      \item \lstinline|struct node| guarda o vizinho \texttt{w} e o ponteiro \texttt{next}.
      \item Espaço proporcional a $V+A$, ou seja, ao \textbf{tamanho} do grafo.
      \item Percorrer os vizinhos de $v$ custa tempo proporcional ao grau de saída de $v$.
      \item \texttt{GRAPHinsertArc()} é cara porque evita arcos repetidos.
    \end{itemize}
  \end{block}
\end{frame}

% =====================================================================
\section{Qual representação escolher?}

\begin{frame}{Matriz {\small versus} listas}
  \begin{center}
  {\small
  \begin{tabular}{lll}
    \toprule
    & \textbf{Matriz de adjacências} & \textbf{Listas de adjacência} \\
    \midrule
    Espaço                        & proporcional a $V^2$ & proporcional a $V+A$ \\
    ``$v$\texttt{-}$w$ é arco?''  & tempo constante      & proporcional ao grau de saída de $v$ \\
    Percorrer vizinhos de $v$     & proporcional a $V$   & proporcional ao grau de saída de $v$ \\
    Inserir arco                  & tempo constante      & proporcional ao número de arcos \\
    Grafos densos                 & apropriada           & sem vantagem \\
    Grafos esparsos               & desperdiça memória   & econômica \\
    \bottomrule
  \end{tabular}}
  \end{center}
  \medskip
  {\small Não existe representação universalmente melhor: a escolha depende da
  densidade do grafo e das operações que o \emph{seu} algoritmo executa com mais
  frequência.}
\end{frame}

\begin{frame}{Entrada de dados: dois formatos de arquivo}
  \begin{columns}[T]
    \begin{column}{0.48\textwidth}
      {\small \textbf{Arquivo de arcos}}\\
      {\scriptsize Linha 1: $V$. Linha 2: $A$. Depois, $A$ linhas com dois
      inteiros em $0..V-1$.}
      \begin{center}\ttfamily\scriptsize
      \begin{tabular}{l}
        7 \\ 6 \\ 0 1 \\ 0 5 \\ 1 0 \\ 1 5 \\ 2 4 \\ 3 1 \\
      \end{tabular}
      \end{center}
    \end{column}
    \begin{column}{0.48\textwidth}
      {\small \textbf{Arquivo de adjacências}}\\
      {\scriptsize Linha 1: $V$. Depois, $V$ linhas: a que começa com $v$ lista
      todos os vizinhos de $v$.}
      \begin{center}\ttfamily\scriptsize
      \begin{tabular}{l}
        7 \\ 0\ \ \ 1 6 \\ 1\ \ \ 0 6 3 \\ 2\ \ \ 4 \\ 3\ \ \ 1 \\ 4\ \ \ 2 \\ 5 \\ 6\ \ \ 0 1 \\
      \end{tabular}
      \end{center}
    \end{column}
  \end{columns}
  {\scriptsize Em ambos os casos a leitura se faz com \texttt{GRAPHinit()} seguido de
  chamadas a \texttt{GRAPHinsertArc()}.}
\end{frame}

\begin{frame}{Para praticar}
  \begin{block}{Exercícios sugeridos (Feofiloff)}
    {\small
    \begin{enumerate}\small
      \item Escreva \texttt{GRAPHindeg()} e \texttt{GRAPHoutdeg()} nas duas representações.
      \item Calcule o vetor booleano \texttt{isSink[]} que identifica os sorvedouros.
      \item Decida se dois vértices são adjacentes; quanto tempo isso consome em cada representação?
      \item Escreva \texttt{GRAPHremoveArc()} para listas de adjacência.
      \item Escreva funções que convertam uma representação na outra.
      \item Escreva \texttt{UGRAPHinsertEdge()} e \texttt{UGRAPHdegrees()} para grafos não-dirigidos.
    \end{enumerate}}
  \end{block}
  {\scriptsize O tutorial em Jupyter que acompanha esta aula traz esqueletos em C
  para os itens 1, 4 e 5.}
\end{frame}

\begin{frame}{Recapitulando a aula}
  \begin{block}{As cinco ideias centrais}
    \begin{itemize}\small
      \item Grafo é um par (vértices, arcos); arco é par \emph{ordenado} $v$\texttt{-}$w$.
      \item Graus resumem a estrutura local; $\sum \text{grau de saída} = A \le V(V-1)$.
      \item Grafos não-dirigidos: pares antiparalelos vistos como arestas, $E = A/2$.
      \item Matriz de adjacências: espaço $V^2$, adjacência em tempo constante.
      \item Listas de adjacência: espaço $V+A$, percurso de vizinhos proporcional ao grau.
    \end{itemize}
  \end{block}
  {\small \textbf{Próxima aula:} caminhos, ciclos e busca em profundidade (DFS) ---
  os primeiros algoritmos que \emph{usam} estas representações.}
\end{frame}

% =====================================================================
\section*{Referências}
\begin{frame}[allowframebreaks]{Referências}
  \scriptsize
  \nocite{*}
  \bibliographystyle{plain}
  \bibliography{../referencias}
\end{frame}

\end{document}
